\documentclass{article}
\usepackage{amsmath}
\usepackage{amssymb}
\usepackage{physics}
\usepackage{bm}

\title{Quantum Mechanics 3}
\author{Kieran Ody}
\date{August 2026}

\begin{document}

\maketitle

\section*{1. Quantum Mechanics}

\subsection*{1(i) de-Broglie Wavelength}

The de-Broglie wavelength of a particle is given by

\begin{equation*}
\lambda=\frac{h}{p}
\end{equation*}

where $h$ is the Planck constant, $\approx 6.626 \times 10^{-34} \, \mathrm{kg\,m^{2}\,s^{-1}}$, and p is the particle's momentum in $\mathrm{kg\,m\,s^{-1}}$.

\vspace{1\baselineskip}

Assuming classical momentum,

\begin{equation*}
p=mv,
\end{equation*}

and substituting gives

\begin{equation*}
\lambda=\frac{h}{mv}
\end{equation*}

The electron is travelling at

\begin{equation*}
v=\frac{c}{137}
\end{equation*}

Using $m_e=9.109 \times 10^{-31}\,\mathrm{kg}$ and $c=2.998 \times 10^8\,\mathrm{m\,s^{-1}}$,

\begin{equation*}
v \approx 2.188\times10^6\,\mathrm{m\,s^{-1}}
\end{equation*}

Hence

\begin{equation*}
\lambda = \frac{6.626\times10^{-34}}{(9.109\times10^{-31})(2.188\times10^6)}
\end{equation*}

giving

\begin{equation*}
\boxed{\lambda\approx3.32\times10^{-10}\,\mathrm{m} = 0.332\,\mathrm{nm}}
\end{equation*}


\subsection*{1(ii) Equal de-Broglie Wavelengths}

It is established that

\begin{equation*}
\lambda=\frac{h}{mv}
\end{equation*}

Since the electron and neutron have equal de-Broglie wavelengths,

\begin{equation*}
\frac{h}{m_ev_e} = \frac{h}{m_nv_n}
\end{equation*}

Therefore,

\begin{equation*}
m_ev_e=m_nv_n,
\end{equation*}

giving

\begin{equation*}
\frac{v_e}{v_n} = \frac{m_n}{m_e}
\end{equation*}

Using

\begin{equation*}
m_e = 9.11 \times 10^{-31} \,\mathrm{kg},
\qquad
m_n = 1.675 \times 10^{-27} \,\mathrm{kg},
\end{equation*}

we obtain

\begin{equation*}
\boxed{\frac{v_e}{v_n} \approx 1839}
\end{equation*}

For classical particles, the kinetic energy is

\begin{equation*}
K=\frac{1}{2}mv^2
\end{equation*}

Thus

\begin{equation*}
\frac{K_e}{K_n} = \frac{m_ev_e^2}{m_nv_n^2}.
\end{equation*}

Using $m_e v_e = m_n v_n$ gives

\begin{equation*}
\boxed{\frac{K_e}{K_n} = \frac{v_e}{v_n} \approx 1839}
\end{equation*}

For equal de-Broglie wavelengths, the electron travels
approximately 1839 times faster than the neutron and has
approximately 1839 times its kinetic energy.


\subsection*{1(iii) Uncertainty Principle and Short-Lived Particles}

The Heisenberg Uncertainty principle relates the uncertainties in
position and momentum according to

\begin{equation*}
\Delta x\Delta p\approx\frac{\hbar}{2}
\end{equation*}

where $\hbar$ is the reduced Planck constant, $= \frac{h}{2\pi} \approx 1.055 \times 10^{-34} \, \mathrm{kg\,m^{2}\,s^{-1}}$,

\vspace{1\baselineskip}

while the energy-time uncertainty relation is

\begin{equation*}
\Delta E\Delta t \approx \frac{\hbar}{2}
\end{equation*}

For a short-lived particle, the momentum and energy uncertainties
can be approximated as

\begin{equation*}
\Delta p \approx mc
\end{equation*}

and

\begin{equation*}
\Delta E \approx mc^2
\end{equation*}

For the $Z$ boson,

\begin{equation*}
mc^2=91.2\,\mathrm{GeV}
\end{equation*}

Using the position-momentum uncertainty relation,

\begin{equation*}
\Delta x \approx \frac{\hbar}{2mc}
\end{equation*}

Multiplying the numerator and denominator by $c$ gives

\begin{equation*}
\Delta x \approx \frac{\hbar c}{2mc^2}
\end{equation*}

Using $\hbar c\approx197.3\,\mathrm{MeV\,fm}$ and $mc^2=91.2\,\mathrm{GeV}=91200\,\mathrm{MeV}$,

\begin{equation*}
\Delta x \approx \frac{197.3}{2(91200)}\,\mathrm{fm},
\end{equation*}

so

\begin{equation*}
\boxed{\Delta x \approx 1.08 \times 10^{-18}\,\mathrm{m}}
\end{equation*}

Similarly, the energy-time uncertainty relation gives

\begin{equation*}
\Delta t \approx \frac{\hbar}{2mc^2}
\end{equation*}

Using $\hbar=6.582\times10^{-25}\,\mathrm{GeV\,s}$,

\begin{equation*}
\Delta t \approx \frac{6.582 \times 10^{-25}}{2(91.2)}\,\mathrm{s},
\end{equation*}

giving

\begin{equation*}
\boxed{\Delta t \approx 3.61 \times 10^{-27}\,\mathrm{s}}
\end{equation*}


\subsection*{1(iv) A Solution of the Wave equation*}

The wave equation* for disturbance $\psi(x,t) = f(x-ct)$ traveling at speed c is:

\begin{equation*}
\frac{\partial^2\psi}{\partial x^2} = \frac{1}{c^2} \frac{\partial^2\psi}{\partial t^2}
\end{equation*}

Propose the form

\begin{equation*}
\psi(x,t) = A\sin{(kx-\omega t)}
\end{equation*}

where $k=\frac{2\pi}{\lambda}$ and $\omega = 2\pi f$, $f=\frac{c}{\lambda}$, and therefore

\begin{equation*}
\omega = kc
\end{equation*}

Thus,

\begin{equation*}
\psi(x,t) = A\sin{(kx-kct)}
\end{equation*}

The second derivative of $\psi$ with respect to $x$ is:

\begin{equation*}
\frac{\partial^2\psi}{\partial x^2} = -Ak^2\sin{(kx-\omega t)} = -k^2\psi
\end{equation*}

The second derivative of $\psi$ with respect to $t$ is:

\begin{equation*}
\frac{\partial^2\psi}{\partial t^2} = -A \omega^2 \sin{(kx-\omega t)} = - \omega^2 \psi = - k^2 c^2 \psi
\end{equation*}

Multiplying by $\frac{1}{c^2}$,

\begin{equation*}
\frac{1}{c^2} \frac{\partial^2\psi}{\partial t^2} = -k^2\psi
\end{equation*}

Thus,

\begin{equation*}
\frac{\partial^2\psi}{\partial x^2} = \frac{1}{c^2} \frac{\partial^2\psi}{\partial t^2}
\end{equation*}

And $\psi(x,t) = A\sin{(kx-\omega t)}$ is a solution of the wave equation*.


\subsection*{1(v) Wavefunction Collapse}

Wavefunction collapse can be illustrated using the double-slit experiment. When electrons are fired individually at a barrier containing two slits, each electron is described by a wavefunction $\psi(x,t)$. The quantity

\begin{equation*}
|\psi(x,t)|^2
\end{equation*}

represents the probability density of detecting the electron at position $x$.

\vspace{1\baselineskip}

When both slits are open and no measurement is made to determine which slit the electron passes through, the wavefunction can be considered as a superposition of the two possible paths:

\begin{equation*}
\psi = \psi_1 + \psi_2
\end{equation*}

The corresponding probability density is therefore

\begin{equation*}
|\psi|^2 = |\psi_1 + \psi_2|^2,
\end{equation*}

which produces an interference pattern on the detector after many electrons have been recorded.

\vspace{1\baselineskip}

However, if a detector is placed at the slits to determine which slit the electron passes through, a definite result is obtained. For example, the electron may be found to have passed through slit 1. The wavefunction is then said to collapse from a superposition of the two possible paths into the state corresponding to the measured outcome:

\begin{equation*}
|\psi\rangle = |\text{slit 1}\rangle + |\text{slit 2}\rangle \quad \longrightarrow \quad |\psi\rangle = |\text{slit 1}\rangle
\end{equation*}

The interference pattern disappears when which-slit information is obtained. This demonstrates the difference between a quantum superposition, in which multiple possible outcomes are represented by the wavefunction, and a measurement, which produces a single definite outcome.

It is important to note that wavefunction collapse does not necessarily require a conscious observer. The measurement can result from a physical interaction between the electron and a detector. More generally, this process is associated with entanglement between the quantum system and its environment and the resulting loss of quantum coherence.

In the standard interpretation of quantum mechanics, wavefunction collapse can be summarised as the transition

\begin{equation*}
\boxed{\text{superposition of possible outcomes} \;\longrightarrow\; \text{definite measured outcome}}
\end{equation*}


\subsection*{1(vi) Classical Kinetic Energy}

Kinetic energy of a particle with mass $m$ and velocity $v$ is often calculated with

\begin{equation*}
E = \frac{1}{2}mv^2
\end{equation*}

As the equation* for momentum is

\begin{equation*}
p = mv,
\end{equation*}

$p$ can be substituted in, such that

\begin{equation*}
E = \frac{1}{2}pv
\end{equation*}

And again, such that

\begin{equation*}
\boxed{E = \frac{p^2}{2m}}
\end{equation*}


\subsection*{1(vii) Atomic Lattice and de-Broglie Wavelength of Neutrons}

For this question, assign the values:

\begin{align*}
d &= 2.15 \times 10^{-10} \,\mathrm{m}\\
\theta &= 30^{\circ} \\
n &= 1
\end{align*}

For a diffraction pattern:

\begin{align*}
n\lambda &= d\sin{\theta} \\
\lambda &= (2.15 \times 10^{-10}) \sin{30^{\circ}} \\
&= 1.075 \times 10^{-10} \,\mathrm{m}
\end{align*}

The de-Broglie wavelength of the neutrons is

\begin{equation*}
\boxed{\lambda = 1.075 \times 10^{-10} \,\mathrm{m}}
\end{equation*}

As established in part (vi),

\begin{equation*}
E = \frac{p^2}{2m}
\end{equation*}

Substituting $p = \frac{h}{\lambda}$,

\begin{equation*}
E = \frac{h^2}{2m\lambda^2}
\end{equation*}

Using

\begin{align*}
h &\approx 6.626 \times 10^{-34} \,\mathrm{kg\,m^{2}\,s^{-1}} \\
m &= m_n \approx 1.675 \times 10^{-27} \,\mathrm{kg} \\
\lambda &= 1.075 \times 10^{-10} \,\mathrm{m}
\end{align*}

gives

\begin{align*}
E &= \frac{(6.626 \times 10^{-34})^2}{2(1.675 \times 10^{-27})(1.075 \times 10^{-10})^2} \\
&\approx 1.134 \times 10^{-20} \,\mathrm{J} \\
&\approx 0.071 \,\mathrm{eV}
\end{align*}

The energy of the neutrons is

\begin{equation*}
\boxed{E = 0.071 \,\mathrm{eV}}
\end{equation*}

The speed of the neutrons is

\begin{align*}
v &= \frac{p}{m_n} \\
&= \frac{h}{m_n\lambda} \\
&= \frac{6.626 \times 10^{-34}}{(1.675 \times 10^{-27})(1.075 \times 10^{-10})} \\
&\approx 3680 \,\mathrm{m\,s^{-1}}
\end{align*}


\subsection*{1(viii) Photoelectric Effect with Silver and UV}

The energy of a photon is

\begin{align*}
E &= hf \\
&= \frac{hc}{\lambda}
\end{align*}

Therefore

\begin{align*}
\lambda &= \frac{hc}{E} \\
&= \frac{(6.626 \times 10^{-34})(2.998 \times 10^8)}{(4.3+5.1) \times (1.602 \times 10^{-19})} \\
&\approx 1.319 \times 10^{-7} \,\mathrm{m}
\end{align*}

The wavelength of the UV light is

\begin{equation*}
\boxed{\lambda = 131.9 \,\mathrm{nm}}
\end{equation*}

An electron of de-Broglie wavelength $1.319 \times 10^{-7}$ m has energy

\begin{align*}
E &= \frac{h^2}{2m_e\lambda^2} \\
&= \frac{(6.626 \times 10^{-34})^2}{2(9.109 \times 10^{-31})(1.319... \times 10^{-7})^2} \\
&\approx 1.385 \times 10^{-23} \,\mathrm{J} \\
&\approx 8.645 \times 10^{-5} \,\mathrm{eV} \\
&\boxed{E = 8.645 \times 10^{-5} \,\mathrm{eV}}
\end{align*}


\subsection*{1(ix) Electron Microscope}

As

\begin{equation*}
\text{resolution} \approx \lambda,
\end{equation*}

\begin{equation*}
\lambda \approx 0.42 \,\mathrm{nm} = 4.2 \times 10^{-10} \,\mathrm{m}
\end{equation*}

An electron of de-Broglie wavelength $4.2 \times 10^{-10}$ m has energy

\begin{align*}
E &= \frac{h^2}{2m_e\lambda^2} \\
&= \frac{(6.626 \times 10^{-34})^2}{2(9.109 \times 10^{-31})(4.2 \times 10^{-10})^2} \\
&\approx 1.366 \times 10^{-18} \,\mathrm{J} \\
&\approx 8.528 \,\mathrm{eV} \\
&\boxed{E = 8.53 \,\mathrm{eV}}
\end{align*}


\subsection*{1(x) The Uncertainty Principle and Beta-decay}

The Uncertainty principle estimates

\begin{equation*}
\Delta E \Delta t \geq \frac{\hbar}{2}
\end{equation*}

Thus,

\begin{align*}
\Delta t &\geq \frac{\hbar}{2 \Delta E} \\
\Delta t &\geq \frac{1.055 \times 10^{-34}}{2(5 \times 10^6)(1.602 \times 10^{-19})} \\
&\boxed{\Delta t \geq 6.586 \times 10^{-23} \,\mathrm{s}}
\end{align*}

The Uncertainty principle also estimates

\begin{equation*}
\Delta x \Delta p \geq \frac{\hbar}{2}
\end{equation*}

Thus,

\begin{align*}
\Delta x &\geq \frac{\hbar}{2 \Delta p} \\
\Delta x &\geq \frac{\hbar}{2 \sqrt{2m_e \Delta E}} \\
\Delta x &\geq \frac{1.055 \times 10^{-34}}{2 \sqrt{2(9.109 \times 10^{-31})(5 \times 10^6)(1.602 \times 10^{-19})}} \\
&\boxed{\Delta x \geq 4.367 \times 10^{-14} \,\mathrm{m}}
\end{align*}


\newpage


\section*{2. The Particle in a One-Dimensional Infinite Potential Well -- Including Task 7 Extension}

The infinite potential well (or ``particle in a box'') is one of the simplest exactly solvable quantum systems. Although idealised, it demonstrates several fundamental concepts of quantum mechanics, including wavefunctions, quantised energy levels and Heisenberg's Uncertainty Principle.

\begin{equation*}
\psi(x,t)
\end{equation*}

describes the wavefunction of a particle of mass $m$ in a `box' of potential energy V.

\vspace{1\baselineskip}

The potential is defined as

\begin{equation*}
V(x)=
\begin{cases}
0, & 0<x<a,\\
\infty, & \text{otherwise}.
\end{cases}
\end{equation*}


\subsection*{2(i) Standing Waves}

The particle is confined to the region $0\leq x\leq a$ by an
infinite potential outside the box. Therefore, the wavefunction
must satisfy the boundary conditions

\begin{equation*}
\psi(0,t)=\psi(a,t)=0
\end{equation*}

At the left-hand boundary, $x=0$, the sine term becomes

\begin{equation*}
\sin(0)=0,
\end{equation*}

so that $\psi(0,t)=0$. At the right-hand boundary, $x=a$,

\begin{equation*}
\sin\left(\frac{n\pi a}{a}\right) = \sin(n\pi) = 0
\end{equation*}

This means the points $x=0$ and $x=a$ are nodes of a standing wave.

\vspace{1\baselineskip}

The condition above is satisfied when $n$ is an integer. The value
$n=0$ would result in $\psi(x,t)=0$ everywhere and therefore cannot
represent a particle, so the allowed values are

\begin{equation*}
\boxed{n=1,2,3,\ldots}
\end{equation*}

Thus, $n$ is a positive integer and represents the quantum number
associated with the standing-wave state.

\begin{equation*}
\cos(\omega t)
\end{equation*}

where $\omega = 2\pi f$, $f=c/\lambda$, describes the wavefunction's sinusoidal time dependence. At a fixed position $x$, the wavefunction therefore oscillates with angular frequency $\omega$:

\begin{equation*}
\psi(x,t)\propto\cos(\omega t)
\end{equation*}

The angular frequency is related to the energy of the state by

\begin{equation*}
E=\hbar\omega
\end{equation*}

\vspace{1\baselineskip}

A suitable form for the wavefunction is therefore

\begin{equation*}
\boxed{\psi(x,t) = A_n\sin\left(\frac{n\pi x}{a}\right)\cos(\omega t)}
\end{equation*}

$A_n$ is the normalisation factor. The sine term describes the
spatial standing-wave pattern, while the cosine term
describes its time dependence.


\subsection*{2(ii) A Solution to the Schrödinger Equation}

The time-independent Schrödinger equation is

\begin{equation*}
-\frac{\hbar^2}{2m} \frac{\partial^2\psi}{\partial x^2} + V\psi = E\psi
\end{equation*}

Inside the box, the potential is zero, and the time-independent Schrödinger equation is therefore

\begin{equation*}
-\frac{\hbar^2}{2m} \frac{\partial^2\psi}{\partial x^2} = E\psi
\end{equation*}

Consider the proposed wavefunction

\begin{equation*}
\psi(x,t) = A_n\sin\left(\frac{n\pi x}{a}\right)\cos(\omega t)
\end{equation*}

Differentiating twice with respect to $x$ gives

\begin{equation*}
\frac{\partial^2\psi}{\partial x^2} = -A_n \left(\frac{n\pi}{a}\right)^2 \sin\left(\frac{n\pi x}{a}\right) \cos(\omega t)
= -\left(\frac{n\pi}{a}\right)^2\psi
\end{equation*}

Substituting this into the Schrödinger equation gives

\begin{equation*}
-\frac{\hbar^2}{2m} \left[-\left(\frac{n\pi}{a}\right)^2\psi\right] = E\psi
\end{equation*}

Hence,

\begin{equation*}
\boxed{E(n,a) = \frac{n^2\pi^2\hbar^2}{2ma^2}}
\end{equation*}

For a proton confined to a box of width

\begin{equation*}
a = 0.47\times10^{-15}\,\mathrm{m},
\end{equation*}

the ground-state energy ($n=1$) is approximately

\begin{equation*}
E_1\approx927\,\mathrm{MeV}
\end{equation*}

Since the energy varies with $n^2$, the general energy is

\begin{equation*}
\boxed{E_n\approx927n^2\,\mathrm{MeV}}
\end{equation*}

The proton rest mass energy is

\begin{equation*}
m_pc^2\approx938\,\mathrm{MeV}
\end{equation*}

Consequently,

\begin{equation*}
\frac{E_1}{m_pc^2} \approx \frac{927}{938} \approx 0.988
\end{equation*}

The ground-state energy is approximately $98.8\%$ of the proton's
rest mass energy, $m_pc^2$. The non-relativistic Schrödinger equation
assumes that the kinetic energy is much smaller than the rest energy,
so that

\begin{equation*}
K_e \approx \frac{p^2}{2m}.
\end{equation*}

However, in this case, the predicted confinement energy is comparable
to $m_pc^2$. This indicates that the proton would have a sufficiently
large momentum for relativistic effects to become important. The non-relativistic particle-in-a-box model no longer provides an accurate description of the proton in such a small confinement region. A relativistic approach would instead be required.


\subsection*{2(iii) Time Dependence and Normalisation}

The time-dependent factor $\cos(\omega t)$ can be replaced by the
complex exponential $e^{-iEt/\hbar}$. This follows from the
relationship between energy and angular frequency,

\begin{equation*}
E=\hbar\omega
\end{equation*}

Thus, the wavefunction can be written as

\begin{equation*}
\psi_n(x,t) = A_n \sin\left(\frac{n\pi x}{a}\right) e^{-iE_nt/\hbar}
\end{equation*}

The Born interpretation states that $|\psi(x,t)|^2$ represents the
probability density of finding the particle at position $x$. Since
the particle must be found somewhere within the box, the total
probability must be total certainty:

\begin{equation*}
\int_0^a |\psi_n(x,t)|^2\,dx=1
\end{equation*}

Taking the modulus squared of the wavefunction gives

\begin{equation*}
|\psi_n(x,t)|^2 = A_n^2 \sin^2\left(\frac{n\pi x}{a}\right) \left|e^{ iE_nt/\hbar}\right|^2
\end{equation*}

Since

\begin{equation*}
\left|e^{-iE_nt/\hbar}\right|^2 = 1,
\end{equation*}

the probability density is independent of time:

\begin{equation*}
|\psi_n(x,t)|^2 = A_n^2 \sin^2\left(\frac{n\pi x}{a}\right)
\end{equation*}

Normalisation therefore requires

\begin{equation*}
A_n^2 \int_0^a \sin^2\left(\frac{n\pi x}{a}\right)\,dx = 1
\end{equation*}

Using

\begin{equation*}
\int_0^a \sin^2\left(\frac{n\pi x}{a}\right)\,dx = \frac{a}{2},
\end{equation*}

we obtain

\begin{equation*}
A_n^2\frac{a}{2}=1
\end{equation*}

Therefore,

\begin{equation*}
\boxed{A_n=\sqrt{\frac{2}{a}}}
\end{equation*}

The normalised wavefunction is

\begin{equation*}
\boxed{\psi_n(x,t) = \sqrt{\frac{2}{a}} \sin\left(\frac{n\pi x}{a}\right) e^{-iE_nt/\hbar}}
\end{equation*}



\subsection*{2(iv) Expected Values of Momentum}

The momentum operator in one dimension is

\begin{equation*}
\hat{p}=-i\hbar\frac{\partial}{\partial x}
\end{equation*}

The expected value of momentum is given by

\begin{equation*}
\langle p\rangle = \int_0^a \psi^*\hat{p}\psi\,dx
\end{equation*}

For the normalised wavefunction

\begin{equation*}
\psi_n(x,t) = \sqrt{\frac{2}{a}} \sin\left(\frac{n\pi x}{a}\right) e^{-iE_nt/\hbar},
\end{equation*}

differentiation with respect to $x$ gives

\begin{equation*}
\frac{\partial\psi_n}{\partial x} = \sqrt{\frac{2}{a}} \frac{n\pi}{a} \cos\left(\frac{n\pi x}{a}\right) e^{-iE_nt/\hbar}
\end{equation*}

Therefore,

\begin{equation*}
\langle p\rangle = -i\hbar\frac{2n\pi}{a^2} \int_0^a \sin\left(\frac{n\pi x} {a}\right) \cos\left(\frac{n\pi x}{a}\right) \,dx
\end{equation*}

The integral evaluates to zero, giving

\begin{equation*}
\boxed{\langle p\rangle=0}
\end{equation*}

Physically, this occurs because the standing wave contains equal
components travelling in opposite directions, corresponding to equal
and opposite momenta.

\vspace{1\baselineskip}

The square of the momentum operator is

\begin{equation*}
\hat{p}^2 = -\hbar^2\frac{\partial^2}{\partial x^2}
\end{equation*}

From the second derivative of the wavefunction,

\begin{equation*}
\frac{\partial^2\psi_n}{\partial x^2} = -\left(\frac{n\pi}{a}\right)^2\psi_n
\end{equation*}

Substituting,

\begin{equation*}
\hat{p}^2\psi_n = \frac{n^2\pi^2\hbar^2}{a^2}\psi_n
\end{equation*}

Taking the expected value gives

\begin{equation*}
\langle p^2\rangle = \frac{n^2\pi^2\hbar^2}{a^2} \int_0^a|\psi_n|^2\,dx
\end{equation*}

Since the wavefunction is normalised,

\begin{equation*}
\int_0^a|\psi_n|^2\,dx=1,
\end{equation*}

and hence

\begin{equation*}
\boxed{ \langle p^2\rangle = \frac{n^2\pi^2\hbar^2}{a^2}}
\end{equation*}


\subsection*{2(v)(a) Expected Values of $x$}

The expected value of x is given by

\begin{equation*}
\langle x\rangle = \int_0^a x|\psi_n|^2\,dx = \frac{2}{a} \int_0^a x\sin^2\left(\frac{n\pi x}{a}\right)\,dx
\end{equation*}

Quoting the integral identity:

\begin{equation*}
\int x \sin^2(\alpha x) \, dx = \frac{1}{4}x^2 - \frac{x\sin(2\alpha x)}{4\alpha} - \frac{\cos(2\alpha x)}{8\alpha^2}
\end{equation*}

Let

\begin{equation*}
\alpha = \frac{n\pi}{a}
\end{equation*}

Evaluating the definite integral from $x = 0$ to $x = a$ (noting that $2\alpha a = 2n\pi$, so $\sin(2n\pi) = 0$ and $\cos(2n\pi) = 1$):

\begin{align}
\int_0^a x \sin^2(\alpha x) \, dx &= \left[ \frac{1}{4}x^2 - \frac{x\sin(2\alpha x)}{4\alpha} - \frac{\cos(2\alpha x)}{8\alpha^2} \right]_0^a \\
&= \left( \frac{1}{4}a^2 - 0 - \frac{1}{8\alpha^2} \right) - \left( 0 - 0 - \frac{1}{8\alpha^2} \right) \\
&= \frac{1}{4}a^2
\end{align}

Thus:
\begin{equation*}
\boxed{\langle x \rangle = \frac{2}{a} \left( \frac{1}{4}a^2 \right) = \frac{1}{2}a}
\end{equation*}


\subsection*{2(v)(b) Calculation of $\langle x^2 \rangle$}

The expected value $\langle x^2 \rangle$ is:

\begin{equation*}
\langle x^2 \rangle = \int_0^a x^2 |\psi_n(x)|^2 \, dx = \frac{2}{a} \int_0^a x^2 \sin^2(\alpha x) \, dx
\end{equation*}

Quoting the integral identity:

\begin{equation*}
\int x^2 \sin^2(\alpha x) \, dx = \frac{1}{6}x^3 - \frac{x \cos(2\alpha x)}{4\alpha^2} - \frac{(2\alpha^2 x^2 - 1)\sin(2\alpha x)}{8\alpha^3}
\end{equation*}

Let

\begin{equation*}
\alpha = \frac{n\pi}{a}
\end{equation*}

Evaluating from $x = 0$ to $x = a$:

\begin{align*}
\int_0^a x^2 \sin^2(\alpha x) \, dx &= \left[ \frac{1}{6}x^3 - \frac{x \cos(2\alpha x)}{4\alpha^2} - \frac{(2\alpha^2 x^2 - 1)\sin(2\alpha x)}{8\alpha^3} \right]_0^a \\
&= \left( \frac{1}{6}a^3 - \frac{a}{4\alpha^2} - 0 \right) - (0 - 0 - 0) \\
&= \frac{a^3}{6} - \frac{a}{4\alpha^2}
\end{align*}

Substituting $\alpha = \frac{n\pi}{a}$:

\begin{equation*}
\int_0^a x^2 \sin^2(\alpha x) \, dx = \frac{a^3}{6} - \frac{a^3}{4n^2\pi^2}
\end{equation*}

And

\begin{align*}
\langle x^2 \rangle &= \frac{2}{a} \left( \frac{a^3}{6} - \frac{a^3}{4n^2\pi^2} \right) \\
&= \frac{a^2}{3} - \frac{a^2}{2n^2\pi^2} \\
&= \frac{1}{3}a^2 \left( 1 - \frac{3}{2n^2\pi^2} \right)
\end{align*}


\subsection*{2(vi)(a) Calculation of $\Delta x$}

The uncertainty in position $\Delta x$ is defined by:

\begin{equation*}
(\Delta x)^2 = \langle x^2 \rangle - \langle x \rangle^2
\end{equation*}

Using the previously derived results of $\langle x \rangle$ and $\langle x^2 \rangle$:

\begin{align*}
(\Delta x)^2 &= \left( \frac{a^2}{3} - \frac{a^2}{2n^2\pi^2} \right) - \left( \frac{a}{2} \right)^2 \\
&= \left( \frac{a^2}{3} - \frac{a^2}{4} \right) - \frac{a^2}{2n^2\pi^2} \\
&= \frac{a^2}{12} - \frac{a^2}{2n^2\pi^2}
\end{align*}

Factoring out $\frac{a^2}{4n^2\pi^2}$:
\begin{align*}
(\Delta x)^2 &= \frac{a^2}{4n^2\pi^2} \left( \frac{4n^2\pi^2}{12} - 2 \right) \\
&= \frac{a^2}{4n^2\pi^2} \left( \frac{n^2\pi^2}{3} - 2 \right)
\end{align*}

Taking the square root yields
\begin{equation*}
\Delta x = \frac{a}{2n\pi} \sqrt{\frac{n^2\pi^2}{3} - 2}
\end{equation*}

\subsection*{2(vi)(b) Calculation of $\Delta p$}

The uncertainty in momentum $\Delta p$ is defined by:
\begin{equation*}
(\Delta p)^2 = \langle p^2 \rangle - \langle p \rangle^2
\end{equation*}

Using the previously derived results of $\langle p \rangle$ and $\langle p^2 \rangle$:
\begin{equation*}
\Delta p = \sqrt{\frac{n^2\pi^2\hbar^2}{a^2} - 0} = \frac{n\pi\hbar}{a}
\end{equation*}

\subsection*{2(vi)(c) Product of Uncertainties $\Delta x \Delta p$}

Multiplying $\Delta x$ and $\Delta p$:
\begin{align*}
\Delta x \Delta p &= \left( \frac{a}{2n\pi} \sqrt{\frac{n^2\pi^2}{3} - 2} \right) \left( \frac{n\pi\hbar}{a} \right) \\
&= \frac{1}{2}\hbar \sqrt{\frac{n^2\pi^2}{3} - 2}
\end{align*}

\subsection*{2(vi)(d) Verification of the Uncertainty Principle}

Heisenberg's Uncertainty Principle requires that:
\begin{equation*}
\Delta x \Delta p \ge \frac{\hbar}{2}
\end{equation*}

For our result to satisfy this inequality, we need:
\begin{equation*}
\sqrt{\frac{n^2\pi^2}{3} - 2} \ge 1
\end{equation*}

Testing for the ground state ($n = 1$):
\begin{equation*}
\frac{(1)^2\pi^2}{3} - 2 \approx 1.2899 > 1
\end{equation*}

Taking the square root:
\begin{equation*}
\sqrt{1.2899} \approx 1.1357 > 1
\end{equation*}

Since $\sqrt{\frac{n^2\pi^2}{3} - 2}$ increases as $n$ increases, the factor is strictly greater than 1 for all quantum numbers $n$, therefore,
\begin{equation*}
\Delta x \Delta p > \frac{\hbar}{2} \quad \text{for all } n = 1, 2, 3, \dots
\end{equation*}

This verifies that the particle in a box always satisfies Heisenberg's Uncertainty Principle.

\end{document}